\section{Cloud Migration Strategies - The 7Rs}
%TODO add diagram
https://aws.amazon.com/blogs/enterprise-strategy/new-possibilities-seven-strategies-to-accelerate-your-application-migration-to-aws/

\begin{itemize}
    \item Retire
    \begin{itemize}
        \item Turn off things you don't need (maybe as a result of Re-architecting)
        \item Helps with reducing the surface for attacks (more security)
        \item Save cost, maybe up to 10\% to 20\%
        \item Focus your attention on resource that must be maintained
    \end{itemize}
    \item Retain
    \begin{itemize}
        \item Do nothing for now (its still a decision to make in a Cloud Migration)
        \item Security, data compliance, performance, unresolved dependencies
        \item No business value to migrate, mainframe or mid-range and non-x86 Unix apps
    \end{itemize}
    \item Relocate
    \begin{itemize}
        \item Move apps from on-premises to its Cloud version
        \item Move EC2 instances to a different VPC, AWS account or AWS region
        \item Example: Transfer servers from VMWare Software-defined Data Centre (SSDC) to VMware Cloud on AWS
    \end{itemize}
    \item Rehost "lift and shift"
    \begin{itemize}
        \item Simple migrations by re-hosting on AWS (applications, databases, data......)
        \item Migrate machines (physical, virtual, another cloud) to AWS cloud
        \item No cloud optimisations being done, applications is migrated as is
        \item Cloud save as much as 30\% on cost
        \item Example: Migrate using AWS application migration service
    \end{itemize}
    \item Replatform "lift and reshape"
    \begin{itemize}
        \item Example: migrate your database to RDS
        \item Example: migrate your application to Elastic Beanstalk
        \item Not changing the core architecture but leverage some Cloud optimisations
        \item Save time and money by moving to a fully manged service or Serverless
    \end{itemize}
    \item Repurchase "drop and shop"
    \begin{itemize}
        \item Moving to a different product while moving to the Cloud
        \item Often you move to SaaS platform
        \item Expensive in the short term, but quick to deploy
        \item Example: CRM to Salesforce.com, HR to Workday, CMS to Drupal
    \end{itemize}
    \item Refactor / Re-architect
    \begin{itemize}
        \item Reimagining how the application is architected using Cloud Native features
        \item Driven by the need of the business to add features and improve scalability, performance, security and agility
        \item Move from a monolithic application to micro-services
        \item Example: move an application to Serverless architectures, use AWS S3
    \end{itemize}
\end{itemize}


\section{Storage Gateway}
\begin{itemize}
    \item Bridge between on-premises data and cloud data
    \item Use cases:
    \begin{itemize}
        \item disaster recovery
        \item backup and restore
        \item tiered storage
        \item on-premises cache and low-latency file access
    \end{itemize}
    \item Types of Storage Gateway
    \begin{itemize}
        \item S3 File Gateway
        \item Volume Gateway
        \item Tape Gateway
    \end{itemize}
\end{itemize}

\subsection{Amazon S3 File Gateway
\begin{itemize}
    \item Configured S3 buckets are accessible using the NFS and SMB protocol
    \item Most recently used data is cached in the file gateway
    \item Supports S3 Standard, S3 Standard IA, S3 One Zone A, S3 Intelligent Tiering
    \item Transition to S3 Glacier using a lifecycle policy
    \item Bucket access using IAM roles for each File Gateway
    \item SMB Protocol has integration with Active Directory (AD) for user authentication
\end{itemize}
%TODO add diagrma

\subsection{Volume Gateway}
\begin{itemize}
    \item Block storage using iSCSI protocol backed by S3
    \item Backed by EBS snapshots which can help restore on-premises volumes!
    \item Cached Volumes: low latency access to most recent data
    \item Stored Volumes: entire dataset is on premise, scheduled backups to S3
\end{itemize}
%TODO add diagram

\subsection{Tape Gateway}
\begin{itemize}
    \item Some companies hav backup processes using physical tapes
    \item With tape gateway, companies use the same processes but, in the cloud
    \item Virtual Tape Library (VTL) backed by S3 and glacier
    \item Back up data using existing tape-based processes (and iSCSI interface)
    \item Works with leading backup software vendors
\end{itemize}
%TODO add diagram

\subsection{AWS Storage Gateway}
%TODO add diagram

\subsection{File Gateway: Extentions}
%TODO add diagrma

\subsection{File Gateway: Read only replicas}
%TODO add diagram

\subsection{File Gateway: Backup and Lifecycle Policies}
%TODO add diagram

\subsection{File Architectures: Other Possibilities}
\begin{itemize}
    \item Amazon S3 ObjectVersioning
    \begin{itemize}
        \item Ability to store multiple object versions as they are modified
        \item Helpful to restore a file to a previous version
        \item Could restore an entire file system to a previous version
        \item Must use the "RefreshCache" API on the gateway to be notified of restore
    \end{itemize}

    \item Amazon S3 object lock
    \begin{itemize}
        \item Enables to have the File Gateway for Write Once Read Many (WORM) data
        \item If there are file modifications or renames in the file share clients, the file gateway creates a new version fo the object without affecting priori verions and the original locked version will remain unchanged
    \end{itemize}
\end{itemize}


%\section{Storage Gateway - Advanced Concepts}


\section{Snow Family}

\subsection{AWS Snowball}
\begin{itemize}
    \item Highly-secure portable devices to collect and process data a tthe edge and migraate data into and out of AWS
    \item Helps migrate up to petabytes of data
\end{itemize}

\begin{itemize}
    \item Device - Snowball Edge Storage Optimised
    \item Compute - 104vCPUs
    \item Memory - 416GB
    \item Storage - 210 TB
\end{itemize}

\begin{itemize}
    \item Device - Snowball Edge Compute Optimised
    \item  Compute - 104vCPUs
    \item Memory - 416GB
    \item Storage (SSD) - 28TB
\end{itemize}

\subsection{Data Migrations with Snowball}

\begin{table}[h]
    \centering
    \begin{tabular}{|c|c|c|c|}
        \hline
        Time to transfer & Column 2 & Column 3 & Column 4 \\
        \hline
        Row 1            & 100MBps  & 1Gbps    & 10Gbps   \\
        \hline
        10 TB            & 12 Days  & 30 hours & 3 hours  \\
        \hline
        100TB            & 124 days & 12 days  & 30 hours \\
        \hline
        1 PB             & 3 years  & 124 days & 12 days  \\
        \hline
        \hline
    \end{tabular}
    \caption{Table caption}
\end{table}

\subsubsection{Challenges}
\begin{itemize}
    \item Limited connectivity
    \item Limited bandwidth
    \item High network cost
    \item Shared bandwidth (can't maximise the line)
    \item Connection stability
\end{itemize}

\subsubsection{AWS Snowball: offline devices to perform data migrations}
- If it takes more than a week to transfer over the network use snowball devices!

\subsection{Diagrams}
- Direct upload to S3:
%TODO add diagrma
- With Snowball
%TODO add diagram

\subsection{What is Edge Computing?
\begin{itemize}
    \item Process data while its being created on an edge location
    --- A truck on the road, a ship on the sea, mining station underground........
    \item The locations may have limited internet and no access to computing power
    \item We setup a snowball edge device to do edge computing
    \begin{itemize}
        \item Snowball edge compute optimised (dedicated for that use case) and Storage Optimised
        \item Run EC2 instances or Lambda function at the edge
    \end{itemize}
    \item Use cases: preprocess data, machine learning, transcoding media
\end{itemize}


\section{Snow Family - Improving Performance}

\begin{itemize}
    \item Most impactful to least:
    \begin{itemize}
        \item Perform multiple write operations at one time - from multiple terminals
        \item Transfer small files in batches - zip up small fields until at least 1Mb
        \item Don't perform other operations on files during transfer
        \item Reduce local network use
        \item Eliminate unnecessary hops - directly connect to the computer
    \end{itemize}

    \item The data transfer rate using the file interface is typically between 25 MB/s and 40MB/s.\\
    If you need to transfer data faster than this use the Amazon S3 Adapter for Snowball, which has a data transfer rate typically between 250Mb/s and 400MB/s
\end{itemize}


\section{AWS DMS - Database Migration Services}
%TODO diagram
- Quickly and securely migrate databases to AWS, resilient, self healing
- The source database remains available during the migration
- Supports:
- Homogenous migrations; ex Oracle to Oracle
- Heterogeneous migrations: Microsoft SQL Server to Aurora
- Continuous Data Replication using CDC
- You must create an EC2 instance to perform the replication tasks

\subsection{DMS Sources and Targets}
\begin{itemize}
    \item Sources
    \begin{itemize}
        \item On-premises and EC2 instances databases: Oracles, MS SQL Server, MySQL, MariaDB, PostgreSQL, MongoDB, SAP, DB2
        \item Azure: Azure SQL Database
        \item Amazon RDS: all including Aurora
        \item Amazon S3
        \item DocumentDB
    \end{itemize}

    \item Targets
    \begin{itemize}
        \item On-premises and EC2 instances databases: Oracle, MS SQL Server, MySQL, MariaDB, PostgreSQL, SAP
        \item Amazon RDS including Aurora
        \item Amazon Redshift
        \item Amazon DynamoDB
        \item Amazon S3
        \item OpenSearch Service
        \item Kinesis Data Streams
        \item DocumentDB
    \end{itemize}
\end{itemize}

\subsection{AWS Schema Conversion Tool (SCT)}
\begin{itemize}
    \item Convert your Databases Schema from one engine to another
    \item Example OLTP: (SQL Server to Oracle) to MySQL, PostgreSQL, Aurora
    \item Example OLAP: (Teradata or Oracle) to Amazon Redshift
%TODO diagram
    \item You do not need to use SCT if you are migrating the same DB engine
    \begin{itemize}
        \item Ex: on-premises PostgreSQL => RDS PostgreSQL
        \item The DB engine is still PostgreSQL (RDS is the platform)
    \end{itemize}
\end{itemize}

\subsection{DMS - Good things to know}
\begin{itemize}
    \item Works over VPC peering, VPN (site to site, software), Direct Connect
    \item Supports Full Load, Full Load + CDC or CDC only
    \item Oracle
    \begin{itemize}
        \item Source: Supports TDE for the source using "BinaryReader"
        \item Target: Supports BLOBs in tables that have a primary key and TDE
    \end{itemize}
    \item OpenSearch:
    \begin{itemize}
        \item Source: does not exist
        \item Target: possible to migrate from a relational database using DMS
        \item Therefore, DMS cannot be used to replica OpenSearch data
    \end{itemize}
\end{itemize}

\subsection{Snowball + Database Migration Service (DMS)}
\begin{itemize}
    \item Larger data migrations can include many terabytes of information
    \item Can be limited due to network bandwidth or size of data
    \item AWS DMS can use Snowball Edge and Amazon S3 to speed up migration
    \item Following stages:
    \begin{enumerate}
        \item You use the AWS Schema Conversion Tool (AWS SCT) to extract the data locally and move it to an edge device
        \item You ship the Edge device or devices back to AWS
        \item After AWS receives your shipment, the edge device automatically loads its data into an Amazon S3 bucket
        \item AWS DMS takes the files an migrates the data t othe target data store. If you are using change data capture (CDC), those updates are written to the Amazon S3 bucket and then applie to the target data store.
    \end{enumerate}
\end{itemize}


\section{AWS CART - Cloud Adoption Readiness Tool}

\subsection{AWS Cloud Adoption Tool (CART)}
\begin{itemize}
    \item Helps organisations develop efficient and effective plans for cloud adoptions and migrations
    \item Transforms your idea of moving to the cloud into a detailed plan that follows AWS best practices
    \item Answer a set of questions across six perspectives (business, people, process, platform, operations, security)
    \item Generates a custom report on your level of migration readiness
\end{itemize}
%TODO see diagram


\section{Disaster Recovery}
\begin{itemize}
    \item Any event that has a negative impact on a companys business continuity or finances is a disaster
    \item Disaster recovering (DR) is about preparing for and recovering from a disaster
    \item What kind of disaster recovery
    \begin{itemize}
        \item On-premises => on-premises: Traditional DR and very expensive
        \item On-premises => AWS cloud: hybrid recovery
        \item AWS Cloud Region A => AWS Cloud Region B
    \end{itemize}
    \item Nee to define two terms:
    \begin{itemize}
        \item Recovery point objective RPO
        \item Recovery Time Objective RTO
    \end{itemize}
\end{itemize}

\subsection{RPO and RTO}
%TODO see diagram

\subsection{Disaster Recovery Strategies}
\begin{itemize}
    \item Backup and Restore
    \item Pilot Light
    \item Warm Standby
    \item Hot Site/ Multi Site Approach
\end{itemize}
%TODO add diagram

\subsection{Backup and Restore (High RPO)}
%TODO add diagram

\subsection{Disaster Recovery - Pilot Light}
\begin{itemize}
    \item A small version of the app is always running in the cloud
    \item Useful for the critical core (pilot light)
    \item Very similar to backup and restore
    \item Faster than Backup and Restore as critical systems are already up
\end{itemize}
%TODO add diagrma

\subsection{Warm Standby}
\begin{itemize}
    \item Full system is up and running, but at minimum size
    \item Upon disaster, we can scale to production load
\end{itemize}
%TODO add diagram

\subsection{Multi Site / Hot Site Approach}
\begin{itemize}
    \item Very low RTO (minutes or seconds) - very expensive
    \item Full production Scale is running AWS and on Premise
\end{itemize}
%TODO add diagram

\subsection{All AWS multi region}
%TODO add diagram

\subsection{Disaster Recovery Tips}
\begin{itemize}
    \item Backup
    \begin{itemize}
        \item EBS Snapshots, RDS automated backups / Snapshots etc
        \item Regulare pushes to S3 / S3 IA / Glacier, Lifecycle Policy, Cross Region Replication
        \item From on-premises, Snowball or Storage Gateway
    \end{itemize}
    \item High Availabilty
    \begin{itemize}
        \item Use Route53 to migrate DNS over from Region to Region
        \item RDS Multi-AZ, ElastiCache Multi-AZ, EFS, S3
        \item Site to Site VPN as a recovery from Direct Connect
    \end{itemize}
    \item Replication
    \begin{itemize}
        \item RDS Replication (Cross Region), AWS Aurora + Global Database
        \item Database Replication from on-premises to RDS
        \item Storage Gateway
    \end{itemize}
    \item Automation
    \begin{itemize}
        \item CloudFormation / Elastic Beanstalk to re-create a whole new environment
        \item Recover / Reboot EC2 instances with CloudWatch if alarms faail
        \item AWS Lambda functions for customised automations
    \end{itemize}
    \item Chaos
    \begin{itemize}
        \item Netflix has a "simian army" randomly terminating EC2
    \end{itemize}
\end{itemize}


\section{AWS FIS - Fault Injection Simulator}
\begin{itemize}
    \item A fully managed service for running fault injection experiments on AWS workloads
    \item Based on Chaos engineering - stressing an application by creating disruptive events (e.g sudden increase in CPU or memeory), observing how the system responds and implementing improvements
    \item Helps you uncover hidden bugs and performance bottlenecks
    \item Supports the following AWS services: EC2, ECS, EKSc, RDS
    \item Use pre-built templates that generate the desired disruptions
\end{itemize}
%TODO add diagram


\section{VM Migrations Serrvices}

\subsection{AWS Application Discovery Service}
\begin{itemize}
    \item Plan migration projects by gathering information about on-premises data centres
    \item Server utilisation data and dependency mapping are important for migraitons
    \item Agentless Discovery (AWS Agentless Discovery Connector)
    \item Open Virtual applicant (OVA) package that can be deplyed to a VMWare host
    \item VM inventory, configuration and performance history such as CPU, memory and disk usage
    \item OS Agnostic
    \item Agent-based Discovery (AWS Application Discovery Agent)
    \item System configuration, system performance, running processes and details of the network connection between systems
    \item Support microsoft server, amazon linux, Ubuntu, RedHat, CentOS, SUSE
    \item Result data can be exported as CSV or viewed within AWS migration hub
    \item Data can be explored using pre-defined queries in Amazon Athena
\end{itemize}

\subsection{AWS Application Discovery Service - Migration Hub Data Exploration}
\begin{itemize}
    \item Allows you to use Amazon Athena to analyse data collected from on-premises servers during discovery
    \item Data is automatically stored in S3 bucket at regular intervals
    \item Use Pre-defined or customer queries in Amazon Athena to analyse data
    \item Example: type of processes running on each server
    \item Ability to upload additional data source such as Configuration Management Database (CMDB) exports
    \item Integrate Athena with QuickSight to visualise data
\end{itemize}
%TODO add diagram

\subsection{AWS Application Migration Service (MGN)}
\begin{itemize}
    \item The "AWS evolution" of CloudEndure Migration, replacing AWS Server Migration Service (SMS)
    \item Lift-and-shift (rehost) solution which simplify migrating application to AWS
    \item Converts your physical, virtual and cloud-based server to run natively on AWS
    \item Supports wide range of platforms, Operating Systems and databases
    \item Minimal downtime, reduced costs
\end{itemize}
%TODO add diagram

\subsection{AWS Elastic Disaster Recovery (DRS)}
\begin{itemize}
    \item Used to be named (CloudEndure Disaster Recovery)
    \item Quickly and easily recover your physical, virtual and cloud-based servers into AS
    \item Example: protect your most critical databases (including Oracle, MySQL and SQL Server)
    \item Continuous block-level replication
\end{itemize}
%TODO diagram

\subsection{On-premises strategy with AWS}
\begin{itemize}
    \item Ability to download Amazon Linux 2 AMI as a VM (.iso format)
    -- VMWare, KVM, VirtualBox (OracleVM), Microsoft Hyper-V
    \item AWS Application Discovery Service
    \begin{itemize}
        \item Gather information about your on-premises servers to plan a migration
        \item Server utilisation and dependency mappings
        \item Track with AWS migration hub
    \end{itemize}
    \item AWS Application Migration Service (MGN)
    \begin{itemize}
        \item Replacing AWS Server Migration Services and CloudEndure Migration
        \item Incremental replication of on-premises live servers to AWS
        \item Migrates the entire VM into AWS
    \end{itemize}
    \item AWS Elastic Disaster Recovery (DRS)
    \begin{itemize}
        \item Replacing CLoudEndure Disaster Recovery
        \item Recover on-premises workloads onto AWS
    \end{itemize}
    \item AWS Database Migration Service (DMS)
    \begin{itemize}
        \item replicate on-premises => AWS, AWS => AWS, AWS => on-premises
        \item Works with various database technologies (oracle, MySQL, DynamoDB, etc)
    \end{itemize}
\end{itemize}


\section{AWS Migration Evaluator}
\begin{itemize}
    \item Helps you build a data-driven business case for migration to AWS
    \item Provides a clear baseline of what you organisation is running to day
    \item Install Agentless Collector to conduct broad-based discovery
    \item Take a snapshot of on-premises foot-print, server dependencies
    \item Analyse current state, define target state then develop migration plan
\end{itemize}
%TODO add diagram


\section{AWS Backup}
\begin{itemize}
    \item Full managed service
    \item Centrally manage and automate backups across AWS services
    \item No need to create custom scripts and manual processes
    \item Supported services
    \begin{itemize}
        \item Amazon EC2 / Amazon EBS
        \item Amazon S3
        \item Amazon RDS (all DBs engines) / Amazon Aurora / Amazon Dyanamo DB
        \item Amazon DocumentDB / Amazon Neptune
        \item Amazon EFS / Amazon FSx (Lustre and Windows File Server)
        \item AWS Storage Gateway (Volume Gateway)
    \end{itemize}
    \item Supports cross-region backups
    \item Support cross-account backups
    \item Supports PITR for supported services
    \item On-Demand and Scheduled backups
    \item Tag-based backup policies
    \item You create backup policies known as Backup plans
    \begin{itemize}
        \item Backup frequency (every 12 hours, daily, weekly, monthly, cron expression)
        \item Backup window
        \item Transition to Cold Storage (Never, Days, Weeks, Months, Years)
        \item Retention Period (Always, Days, Weeks, Months, Years)
    \end{itemize}
\end{itemize}

%TODO add diagram

\subsection{AWS Backup Vault Lock}
\begin{itemize}
    \item Enforce a WORM (Write Once Read Many) state for all the backups that you store in your AWS backup vault
    \item Additional layer of defence to protect your backups against
    \begin{itemize}
        \item Inadvertent or malicious delete operations
        \item Updates that shorten or alter retention periods
    \end{itemize}
    \item Even the root user cannot delete backups when enabled
\end{itemize}
%TODO add diagrma
